%! Author = Stefan
%! Date = 12.10.2026

\documentclass[a4paper, 11pt]{article}

\usepackage[ngerman]{babel}
\usepackage[margin=2.5cm]{geometry}

\title{Implementierungsanleitung TwoSum}
\author{Stefan Ertl}
\date{}

\begin{document}

    \maketitle

    \section{Ziehle dieses Algorithmus}
    Diese Implementation hat zum Ziel möglichst intuitiv und simpel zu sein.

% Text

    \section{Implementierung des Algorithmus}
    Der Algorithmus erhält als Übergabeparameter zwei Arrays: \textit{Numbers} und \textit{Targets}.\\\\
    Als Erstes werden alle möglichen Summen der Zahlen in \textit{Numbers} gebildet und in einem Array \textit{Sums} zwischengespeichert.\\\\
    Anschließend wird jede Zahl in \textit{Targets} in \textit{Sums} gesucht; ist diese vorhanden, so kann diese in die Ergebnisliste eingetragen werden.
% Text

    \section{Weitere Optimierungen}


% Text

    \section{Laufzeit}
    Für alle n Zahlen in \textit{Numbers} werden erneut alle n Zahlen in \textit{Numbers} durchlaufen. Zusätzlich muss für die Suche für alle m Zahlen in \textit{Targets} alle \(n^2\) Zahlen in \textit{Sums} durchsuchen.\\\\
    Zusammen: \[
                  \mathcal{O}\left(n^2 + m \cdot n^2\right)
    \]

% Text

\end{document}